\section{随机性提取器}\label{sec:extractor}

本章探讨随机性提取器的构造与分析。在 \ref{sec:extractor_special} 中，我们构造了一个具有特殊性质的随机性提取器，这完成了此前遗留的 \cref{lem:extractor_unif} 的证明。在 \ref{sec:lower_multi_ext} 中，作为尝试分析乘法误差的自然想法，我们引入了带有乘法误差的随机性提取器。但遗憾的是，我们反而证明了构造带有乘法误差的随机性提取器是不可能的。

\subsection{带有特殊性质的提取器}\label{sec:extractor_special}

我们在此重述随机性提取器的定义与 \cref{lem:extractor_unif}，并给出其主要构造思路。详细证明见 \ref{sec:omitted_proof_extractor}。与以往工作不同，该提取器额外满足如下性质：对任意固定种子 $s$，均有 $\Ext(U_n, s) = U_m$，称之为\emph{均匀性性质}。

\Extractor*

\ExtractorUniform*

我们的构造总体上遵循 Guruswami、Umans 和 Vadhan \cite{GUV_extractor} 的框架，但在若干关键细节上有所调整。构造中同样需要一个“剩余哈希引理”提取器以及一个凝结器（见 \cref{def:condenser}）。我们采用 Cheraghchi 博士论文 \cite{Cheraghchi11_phd_thesis} 中的构造，其关键优势在于：对任意固定种子 $s$，提取器 $\Ext(\cdot, s)$ 与凝结器 $\Con(\cdot, s)$ 均为 $\F_2$ 上的线性函数。

利用这一线性结构，可以强制获得均匀性性质。具体而言，若 $\Ext(x, s) = M_s \cdot x$ 且 $M_s$ 不满秩，则可将其替换为一个满秩矩阵，从而得到一个满射线性映射。该修改不会降低输出的最小熵，并保证 $\Ext(U_n, s) = U_m$。

在此基础上，可以如 \cite{GUV_extractor} 中一样，通过凝结器与提取器的组合，构造出种子长度与熵同量级的提取器。由于该构造对随机源 $X$ 进行分块，并在不同块上分别应用子提取器，因此均匀性性质可以自然继承。

最后，我们采用 \cite{GUV_extractor} 中的递归式构造进一步降低种子长度。然而，该递归形式
$$
\Ext(x, (s_1, s_2)) = (\Ext_1(x, s_1), \Ext_2(x, s_2))
$$
与分块构造不同，两个子提取器作用于同一个随机源 $X$。即使 $X$ 为均匀分布，也无法保证两部分输出的独立性，从而可能破坏均匀性性质。

为解决这一问题，我们利用提取器的线性结构。设 $\Ext_1(x, s_1) = M_1 \cdot x$，$\Ext_2(x, s_2) = M_2 \cdot x$，且 $M_1, M_2$ 均为满秩矩阵。若 $\begin{pmatrix} M_1 \\ M_2 \end{pmatrix}$ 为双射，则均匀性性质仍然成立。其中一个典型情形是 $M_1$ 与 $M_2$ 互为\emph{正交补}，即满足 $M_1^\top M_2 = 0$。

基于这一观察，在应用 $\Ext_1$ 提取 $X$ 后，不再直接对 $X$ 进行第二次提取，而是考虑其在正交补上的投影 $\overline{X} := M_1^\bot X$。随后，$\Ext_2$ 作用于 $\overline{X}$ 而非 $X$。这样一来，$\Ext_1$ 与 $\Ext_2$ 实际上作用于 $X$ 的两个“正交部分”（尽管这种划分并非显式的分块结构），从而保持均匀性性质。

此外，可以证明，在给定 $M_1 X$ 的条件下，$\overline{X}$ 的最小熵恰好等于 $X$ 的剩余熵，因此 $\overline{X}$ 仍具有足够的随机性以支持进一步提取。

\subsection{乘法误差提取器的下界分析}\label{sec:lower_multi_ext}

最小哈希族的核心要求是目标概率满足乘法误差界。回顾 \cref{thm:min_wise_log_seed} 与 \cref{thm:k_min_wise_klog_seed} 的证明，整体概率被分解为多个子桶事件的连乘积。在误差分析中，连乘积中每项的加性误差在转换为乘法误差时通常都要付出很大代价；若生成子桶种子的提取器本身能提供乘法误差保证，则连乘积的误差传播将天然保持相对误差形式，从而大幅简化分析并有望突破现有参数瓶颈。受此直观动机驱动，我们自然提出一个问题：是否存在显式的非平凡乘法提取器？

在标准设定中，随机性提取器是通过加法误差 $\varepsilon$ 来定义的。针对这一经典定义，Radhakrishnan 和 Ta-Shma \cite{RT_siam00} 建立了关于种子长度和熵损失的一个基本下界，而 Guruswami-Umans-Vadhan 等里程碑式构造 \cite{GUV_extractor} 已基本匹配该下界，证明了显式非平凡提取器的存在性。然而，我们将证明，一旦将误差度量转换为乘法形式，可行参数空间将发生根本性坍缩。为严格表述并建立下界，我们首先给出乘法提取器的形式化定义。

\begin{definition}[乘法提取器]
    设 $\Ext : \{0,1\}^n \times \{0,1\}^d \to \{0,1\}^m$。若对于 $\{0,1\}^n$ 上任意满足最小熵 $\Hmin(X)\ge k$ 的分布 $X$，以及任意子集 $S\subseteq\{0,1\}^m$，均有
    $$
    \Pr[\Ext(X, U_d) \in S] = \Pr[U_m \in S] \cdot (1 \pm \delta),
    $$
    则称 $\Ext$ 为一个 $(k,\delta)$-乘法提取器。
\end{definition}

给定一个 $\Ext$，我们考虑与之关联的提取器图 \cite{RT_siam00}。

\begin{definition}[乘法提取器图]
    给定 $\Ext : \{0, 1\}^n \times \{0, 1\}^d \to \{0, 1\}^m$，令 $N = 2^n, D = 2^d$ 且 $M = 2^m$。其乘法提取器图定义为二分图 $G = (V_1 = [N], V_2 = [M], E)$，满足：对于每一组满足 $\Ext(u, s) = v$ 的 $u \in [N], v \in [M]$ 和 $s \in [D]$，我们将边 $(u, v)$ 加入 $G$。
    
    若对于任意满足 $|X| = K$ 的 $X \subseteq V_1$ 以及任意 $Y \subseteq V_2$，均有
    $$
    \left| \frac{E(X,Y)}{E(X,V_2)}-\frac{|Y|}{|V_2|} \right| \le \delta \cdot \frac{|Y|}{|V_2|},
    $$
    其中 $E(X, Y)$ 表示 $X$ 与 $Y$ 之间在边集 $E$ 中的边数，则称 $G$ 为一个 $(K, \delta)$-乘法提取器图。
\end{definition}

这两个定义是等价的。

\begin{proposition}\label{prop:eql_ext_graph}
    $\Ext$ 是一个 $(k, \delta)$-乘法提取器，当且仅当其对应的图是一个 $(2^k, \delta)$-乘法提取器图。
\end{proposition}

基于上述定义，我们证明：当 $\delta < 1$ 时，任何有效的左正则乘法提取器图都要求其左度数 $D$ 大得平凡（即 $D > M / 2$）。这说明了构造带有乘法误差提取器的不可能性，也凸显了前文构造方法的必要性与创新性。

\begin{theorem}[乘法误差 $\delta < 1$ 提取器的下界]
    设 $G = (V_1 = [N], V_2 = [M], E)$ 是一个左度数为 $D$ 的左正则 $(K, \delta)$-乘法提取器图。若 $\delta < 1$，则
    $$
    D > M \left( 1 - \frac{K}{N} \right).
    $$
    
    特别地，对于满足 $K \le N/2$ 的，也就是最小熵至多是 $k - 1$ 的典型弱随机源，必须有 $D > M/2$，也就是种子长度只能平凡：$d > m - 1$。
\end{theorem}

\begin{proof}
    由于 $G$ 是度数为 $D$ 的左正则图，对于任意满足 $|X| = K$ 的子集 $X \subseteq V_1$，从 $X$ 引出的边总数恰好为 $E(X, V_2) = K \cdot D$。
    
    根据乘法提取器图的定义，对于任意大小为 $K$ 的 $X \subseteq V_1$ 和任意 $Y \subseteq V_2$，我们有
    $$
    \left| \frac{E(X,Y)}{K D} - \frac{|Y|}{M} \right| \le \delta \cdot \frac{|Y|}{M}
    $$
    
    考虑 $V_2$ 中任意顶点 $y$ 构成的单点集 $Y = \{y\}$。该不等式变为
    $$
    \left| \frac{E(X,\{y\})}{K D} - \frac{1}{M} \right| \le \frac{\delta}{M} \quad \Longrightarrow \quad \frac{E(X,\{y\})}{K D} \ge \frac{1 - \delta}{M}.
    $$
    
    由于 $\delta < 1$，故 $\frac{1 - \delta}{M} > 0$。因此，它严格要求 $E(X, \{y\}) > 0$。这意味着\emph{每一个}大小为 $K$ 的子集 $X \subseteq V_1$ 都必须包含至少一个与 $y$ 相邻的顶点。 
    
    换言之，$y$ 在 $V_1$ 中的非邻居集合（记为 $V_1 \setminus \Gamma(y)$）的大小不能达到或超过 $K$。因此，$y$ 在 $V_1$ 中的非邻居数量至多为 $K - 1$。
    
    由于 $|V_1| = N$，$y$ 的度数必须满足
    $$
    d(y) \ge N - (K - 1) = N - K + 1.
    $$
    
    关键在于，这一结论必须对 $V_2$ 中的每一个顶点 $y$ 都成立。我们可以对 $V_2$ 中所有顶点的度数求和，以计算 $G$ 中的总边数：
    $$
    N \cdot D = |E| = \sum_{y \in V_2} d(y) \ge M \cdot (N - K + 1) \quad \Longrightarrow \quad D \ge M \left( 1 - \frac{K - 1}{N} \right) > M \left( 1 - \frac{K}{N} \right).
    $$
    
    这表明当 $K \le N / 2$ 时，种子长度 $D$ 被平凡地限制为至少 $M / 2$，从而使得在此参数范围内构造非平凡提取器成为不可能。
\end{proof}

\begin{remark}[关于 $\delta \ge 1$ 的情形]
    上述不可能性证明严重依赖于 $\delta < 1$ 这一条件。若我们将误差参数放宽至 $\delta \ge 1$，则即使 $E(X, \{y\}) = 0$，该下界约束也会平凡地成立。在此参数范围内，乘法提取器实际上等效于一个零误差凝结器（zero-error condenser）。
    
    为说明这一等价性，考虑令 $\delta = 2^{m-k'} - 1 \gg 1$。乘法误差条件仅对输出分布中任意单点 $y \in \{0,1\}^m$ 的概率施加了上界限制：
    $$
    \Pr[\Ext(X, U_d) = y] \le (1 + \delta) \frac{1}{M} = 2^{-k'}.
    $$
    
    这等价于要求输出分布的最小熵满足 $H_\infty(\Ext(X, U_d)) \ge k'$。这与凝结器的定义完全一致。关键区别在于：传统凝结器通常允许较小的加法误差（即输出分布与某个最小熵为 $k'$ 的分布统计距离不超过 $\varepsilon$），而此处的乘法定义要求该最小熵保证严格成立、不带任何加法误差，因此称之为“零误差”凝结器。
    
    在此 $\delta \ge 1$ 的范围内，种子长度是否存在非平凡的下界，或者此类零误差凝结器的显式构造是否真正可行，仍是值得研究的有趣开放问题。
\end{remark}

\subsection{总结}

本章围绕随机性提取器的构造与其在最小哈希分析中的潜在应用展开研究。我们首先构造了一类具有额外均匀性性质的随机性提取器，从而完成了此前遗留的 \cref{lem:extractor_unif} 的证明。该构造大致遵循 Guruswami、Umans 和 Vadhan \cite{GUV_extractor} 给出的框架，并利用线性性配合正交投影，使得在递归构造的时候可以保持均匀性性质。

进一步地，我们从分析乘法误差传播的角度出发，引入了乘法误差版本的随机性提取器概念，并对其可行性进行了系统研究。我们证明，在标准模型下，非平凡意义上的乘法误差提取器并不存在，从而表明该方向无法直接用于改进最小哈希族的误差分析。

综上，本章一方面提供了满足额外结构性质的提取器构造，另一方面揭示了乘法误差提取器在一般情形下的不可实现性。这一对比结果说明，在最小哈希族的分析中，必须依赖更精细的结构性分解方法，而不能简单诉诸误差形式的直接强化。